\section{Evidence Required for a Brain-Specific Training Advantage}
\label{sec:valid-claim}

Brain-encoding \autocite{oota2023,merlin2024}, brain-response-guided optimization \autocite{moussa2025,negi2025,merlinData2026,vaidya2026}, privileged-information and feature-distillation \autocite{lopezpaz2016,saadi2025}, and model-compression \autocite{oota2026} studies address different questions.
Predicting recorded brain responses does not establish that training with those responses changes a student model.
A change in the student also does not establish that it arises specifically from brain-response content, transfers to independently recorded \ac{fmri} responses, or improves a useful endpoint at comparable \ac{lm} quality.
This section connects these claims and defines the evidence required to support a brain-specific training advantage.

\subsection{What controlled brain predictivity establishes}
\label{sec:controlled-score}

Controlled brain predictivity tests whether representations extracted from a frozen \ac{lm} contain information useful for predicting recorded \ac{fmri} responses to held-out stimuli.
An encoding model is fitted to predict responses from those representations on one set of stimuli and evaluated on a separate set.
This held-out evaluation tests whether prediction generalizes beyond the fitting stimuli rather than merely reflecting overfitting to them \autocite{oota2023,merlin2024}.
Under a linear readout, a positive controlled score means that a linear combination of \ac{lm} representation dimensions improves prediction beyond the declared controls.

Three controls are required to determine whether predictive performance is associated with learned \ac{lm} representations rather than simpler alternatives.
First, a nuisance-only model predicts recorded \ac{fmri} responses from stimulus properties such as position, length, and static lexical features.
Comparing it with a full model that also includes the tested representations measures what those representations add.
Second, an architecture-matched untrained \ac{lm} retains the tokenizer, dimensionality, and architecture of the trained model but lacks its learned parameters.
This comparison tests whether predictive performance depends on learned representations rather than the model's fixed structure.
Third, contiguous or blocked splits keep locally related samples from being scattered across fitting and evaluation sets, reducing inflated scores caused by dependence or leakage between them \autocite{hadidi2026}.
Together, these controls narrow the observed effect to held-out predictive value added by learned \ac{lm} representations under the declared assay.

Even with these controls, controlled brain predictivity remains bounded by the declared assay.
A linear readout tests whether recorded responses can be predicted from a linear combination of \ac{lm} representation dimensions; it does not test every form of information that might be present.
Participant-specific responses measure predictivity for an individual, whereas averaged responses measure predictivity of a group-level response summary.
The stimulus distribution, response definition and aggregation, response reliability, and selected brain regions further determine what can be predicted \autocite{lage2019}.
The inference unit determines the scope of uncertainty: many voxels or repeated folds from one participant do not provide evidence about many independent participants.
Controlled brain predictivity therefore supports a bounded claim about held-out prediction under the specified assay, not a shared biological mechanism, causal equivalence, or general cognitive competence \autocite{jia2026}.

\subsection{Why brain-response targets might help a student}
\label{sec:target-rationale}

Ordinary \ac{kd} trains a student to reproduce the teacher's output distribution, but this objective does not uniquely determine how the student organizes information in its middle layers.
Different bases, subspaces, and feature organizations can support similar token predictions.
Two students can therefore show similar language behavior while retaining different information in their model representations.
This leaves room for an auxiliary training target to favor one representation over another without necessarily improving or degrading the language objective.
Feature-distillation research further shows that the value of such guidance depends on which information is emphasized and how teacher and student capacities are matched, rather than simply on copying a large dense vector \autocite{saadi2025}.

Brain-response targets can exploit this flexibility as privileged information available during training but unnecessary at deployment \autocite{lopezpaz2016}.
The student retains the same \ac{lm} objective, while an auxiliary objective favors representations from which the training target can be predicted.
Recorded \ac{fmri} responses provide observed biological measurements, whereas synthetic brain responses generated from text package structure learned by a fixed generator.
Because each synthetic response is determined by the text and generator, it cannot provide example-specific information that varies independently of them.
It may nevertheless present generator-learned structure in a form that is easier for the student to learn or recover \autocite{cover2006}.

Evidence that the auxiliary objective affected training does not establish a brain-specific benefit.
Reduced target loss may show only that the model learned to predict the auxiliary target during training, without demonstrating a lasting change in the student.
Target uptake instead requires a new prediction model fitted after training to recover the target from frozen retained-student representations relative to the declared baseline.
Retained-student movement separately requires parameter or representation changes that remain in the student after training.
These checks show, respectively, that target information became recoverable and that the intervention changed the retained student.
They do not establish that the changes arose specifically from brain-response content, transferred to independently recorded \ac{fmri} responses, or improved a practical endpoint.

\subsection{What prior brain-response-guided training leaves unresolved}
\label{sec:prior-gap}

Two adjacent literatures bear on the governing question, but they assign neural data different roles.
Post-hoc encoding studies use brain responses to evaluate representations after a model has been changed for another reason.
Narrative-summarization tuning can raise \ac{fmri} encoding scores without using neural supervision, and scaling can raise naturalistic brain alignment while matched-size instruction tuning does not under the studied assay \autocite{aw2023,gao2024}.
These results show that an alignment score responds to non-neural objective and capacity choices; they do not estimate the value of neural data as a training target.
Neural-guided studies instead place recorded or predicted neural responses inside the optimization loop.
Their positive findings establish that such targets can change representations, encoding scores, or selected task endpoints, but the estimand depends on the model intervention, comparator, transfer boundary, and inference unit.
Table~\ref{tab:related-designs} compares designs along those axes rather than treating every positive brain score as the same claim.

\begin{table}[H]
  \centering
  \footnotesize
  \setlength{\tabcolsep}{3pt}
  \newcommand{\relatedarraystretch}{1.08}
  \renewcommand{\arraystretch}{\relatedarraystretch}
  \caption{Design comparison for work closest to the governing question. ``Neural data role'' distinguishes post-hoc evaluation from a training target. The final column states the remaining gap, not a defect in the cited study's own claim.}
  \label{tab:related-designs}
  \begin{tabular}{@{}p{24mm} p{38mm} p{48mm} p{42mm}@{}}
    \toprule
    Design & Neural data role and model change & Strongest relevant comparison or endpoint & Gap relative to this study's question \\
    \midrule
    \textcite{aw2023}; \textcite{gao2024} &
    Post-hoc evaluation after summarization tuning, scaling, or instruction tuning. &
    Base--tuned or matched-size model comparisons show that non-neural training choices can change measured alignment. &
    No neural training target, compressed student, or neural-versus-nonneural intervention contrast. \\
    \textcite{bilgin2026} &
    Participant-\ac{fmri} cosine loss plus token loss updates upper-layer adapters in pretrained text models. &
    Held-out-season and cross-stimulus encoding, new-participant readouts, and a model-specific external-task gain. &
    No shuffled or matched nonbrain target, smaller student, or ordinary-distillation comparator. \\
    \textcite{merlinData2026} &
    Participant-\ac{fmri}, stimulus-language, or joint objectives update adapters in BERT and GPT-2. &
    Same stimulus text and broad tuning budget; brain tuning wins more linguistic-probe comparisons than stimulus tuning. &
    Target geometry and difficulty are unmatched; no participant-held-out transfer or compression comparison. \\
    \textcite{merlin2026} &
    Gradient reversal removes brain-predictive information; complementary tuning increases it. &
    A permuted-\ac{fmri} control and matched language-model loss accompany changes in broad linguistic probes. &
    No compressed student; the manipulated readout may contain linguistic variance not specific to neural content. \\
    \textcite{freteault2025} &
    Individual or group \ac{fmri} fine-tunes selected layers of an auditory network. &
    Held-out-season brain encoding and multiple external auditory tasks improve. &
    No matched nonbrain or shuffled-target fine-tuning arm; the model and modality differ from text distillation. \\
    \textcite{pirlot2022} &
    A neural regularizer aligns a vision network with monkey V1 during classification training. &
    Shuffled neural labels reproduce the clean-accuracy gain, whereas moderate-attack robustness favors real targets. &
    The clean-task gain is generic regularization; only the robustness endpoint separates neural content in that design. \\
    \bottomrule
  \end{tabular}
\end{table}

The broader positive literature spans language, speech, and vision rather than one common treatment effect.
\ac{fmri}-guided speech adaptation has improved encoding and selected semantic or auditory tasks, including transfer across data or participant settings \autocite{moussa2025,moussa2025b,freteault2025}.
Text-model interventions report multilingual, linguistic-probe, visual-commonsense, or reasoning changes after neural guidance \autocite{negi2025,bilgin2026,merlinData2026,xiao2026}.
Temporally precise electrocorticography targets have improved neural prediction within a speech substrate, while slow-\ac{fmri} tuning has transferred to fast electrocorticography prediction for new participants and stimuli after fitting the required readout \autocite{zhang2026,vaidya2026}.
These results make neural-guided optimization empirically plausible.
They do not imply one shared endpoint: some adapt a full encoder or adapters rather than a smaller student, some fit participant-specific targets, and some transfer across stimuli, datasets, participants, or modalities with different amounts of target-participant data.

Controls also answer different alternatives.
A structure-preserving permutation asks whether correct stimulus--target pairing matters.
A task-only or stimulus-language arm asks whether additional exposure or a conventional objective suffices.
A held-out participant or dataset tests a transfer boundary, often after fitting a new readout.
None automatically supplies a nonbrain target matched on dimensionality, covariance, baseline predictability, learnability, loss geometry, and optimization pressure.
The shuffled-label result in vision is especially instructive: a clean-accuracy gain can survive destruction of item-specific neural content even when another endpoint distinguishes the real target \autocite{pirlot2022}.
Positive neural-guidance results should therefore be retained as positive within their designs, while attribution is limited to the controls each design actually includes.

Privileged-information and feature-distillation work supplies the generic-target alternative directly.
Dense training-only supervision can help by exposing a useful subspace or easing optimization, and its value depends on teacher information, student capacity, and the chosen representation rather than the target's scientific origin \autocite{lopezpaz2016,saadi2025}.
Measurement-validity work separately requires blocked splits, nuisance controls, and honest inference units before an encoding score is interpreted \autocite{hadidi2026,lage2019,jia2026}.
Compression studies that measure alignment after pruning or quantization show that shrinkage method and scale can change the score, but they do not test neural supervision during compression \autocite{oota2026}.

The unresolved pre-results question is therefore narrower than whether neural-guided training can ever change a model.
For a fixed smaller student already trained by ordinary \ac{kd}, does a recorded or predicted neural target add value beyond matched non-neural training at comparable language-model quality?
Answering it requires target uptake and retained-student movement, correct-pairing and matched-target controls, participant-level biological transfer, and a separate external endpoint before a brain-specific utility claim is warranted.

\subsection{Evidence standard for a brain-specific training advantage}
\label{sec:evidence-standard}

Figure~\ref{fig:evidence-chain} orders the distinct claims that must be supported, while marking exact-target measurability as an additional prerequisite for predicted rather than recorded neural targets.

\begin{figure}[H]
  \centering
  \begin{tikzpicture}[
    stage/.style={draw, rounded corners, align=center, text width=33mm, minimum height=13mm, inner sep=4pt, font=\small},
    link/.style={-{Latex[length=2mm]}, thick}
  ]
    \node[stage] (measure) {\textbf{1. Controlled measurability}\\[-1mm]\footnotesize Assay accesses recorded responses};
    \node[stage, right=3mm of measure] (headroom) {\textbf{2. Quality-aware headroom}\\[-1mm]\footnotesize A fair intervention opportunity exists};
    \node[stage, right=3mm of headroom] (target) {\textbf{3. Exact-target measurability}\\[-1mm]\footnotesize Predicted-target path only: target relates to recorded responses};
    \node[stage, right=3mm of target] (uptake) {\textbf{4. Target uptake}\\[-1mm]\footnotesize Fresh held-out recovery improves beyond baseline};
    \node[stage, below=5mm of uptake] (movement) {\textbf{5. Retained-student movement}\\[-1mm]\footnotesize Change survives head removal and preserves quality};
    \node[stage, left=3mm of movement] (attribute) {\textbf{6. Comparator validity}\\[-1mm]\footnotesize Attribution to neural target content becomes testable};
    \node[stage, left=3mm of attribute] (transfer) {\textbf{7. Biological transfer}\\[-1mm]\footnotesize Independent recorded responses improve};
    \node[stage, left=3mm of transfer] (utility) {\textbf{8. External utility}\\[-1mm]\footnotesize A practical endpoint improves};
    \draw[link] (measure) -- (headroom);
    \draw[link] (headroom) -- (target);
    \draw[link] (target) -- (uptake);
    \draw[link,dashed] (headroom.north) to[bend left=28] node[above,font=\scriptsize,fill=white,inner sep=1pt] {recorded-target bypass} (uptake.north);
    \draw[link] (uptake) -- (movement);
    \draw[link] (movement) -- (attribute);
    \draw[link] (attribute) -- (transfer);
    \draw[link,dashed] (transfer) -- (utility);
  \end{tikzpicture}
  \caption{Conceptual evidence chain. Each stage states the additional claim licensed by passing it. Section~\ref{sec:framework} later maps the stages to estimands and experiments; the Results report their verdicts.}
  \label{fig:evidence-chain}
\end{figure}

No single comparator can answer every stage.
An architecture-matched untrained network tests whether controlled measurability depends on learned parameters, and an ordinary-distillation arm tests the incremental effect of adding a neural objective to the same student-training regime.
Comparable language-model quality prevents a change in generic model quality from silently deciding an alignment contrast.
A permuted target tests whether correct stimulus--target pairing matters while largely preserving the target's marginal structure.
An adequately matched nonbrain target makes attribution to neural target content testable by controlling geometry, scale, baseline predictability and learnability, and optimization pressure; attribution additionally requires the observed neural-target contrast to favor that arm.
Finally, a freshly fitted readout on independent recorded responses tests biological transfer; success on the optimized target itself cannot substitute for that endpoint.

The chain through biological transfer is conjunctive for a brain-specific biological advantage; external utility is an additional requirement only for a practical-utility claim.
When a prerequisite fails, later target uptake or retained-student movement may still describe what the intervention did, yet it cannot restore the stronger attribution or transfer claim.
An unresolved comparison should remain unresolved rather than being converted into a null, and a failed test under one readout, cohort, or target should not be universalized beyond its scope.
This stopping rule keeps a large downstream score from compensating for an unidentified mechanism earlier in the chain and makes clear which redesign would be needed next.
Section~\ref{sec:framework} maps each gate to its data and training arms, estimand, comparator, endpoint, and inference unit.
